% Generado automáticamente: índice de fuentes con saltos a sección exacta.
% NO editar a mano (regenerar con el script).

\subsection*{\hyperref[src:clase6]{Clase 6 — Políticas y Control de Acceso}}
\begin{itemize}[leftmargin=1.4em,itemsep=0pt,topsep=2pt]
  \item[]{\bfseries \hyperlink{clase6.4}{1~Definiendo seguridad en aplicaciones}}
  \item[]\hspace{1.2em}{\hyperlink{clase6.4}{1.1~La tríada CIA}}
  \item[]{\bfseries \hyperlink{clase6.5}{2~Conceptos de control de acceso}}
  \item[]\hspace{1.2em}{\hyperlink{clase6.6}{2.1~La tríada desde el punto de vista de los accesos}}
  \item[]\hspace{1.2em}{\hyperlink{clase6.6}{2.2~Identidad del sujeto}}
  \item[]{\bfseries \hyperlink{clase6.6}{3~Políticas vs.\ Mecanismos}}
  \item[]{\bfseries \hyperlink{clase6.7}{4~Modelos de seguridad}}
  \item[]{\bfseries \hyperlink{clase6.7}{5~Modelo Bell--LaPadula (confidencialidad)}}
  \item[]\hspace{1.2em}{\hyperlink{clase6.7}{5.1~Reglas (principios)}}
  \item[]\hspace{1.2em}{\hyperlink{clase6.9}{5.2~Compartimentos, reticulado y dominancia}}
  \item[]{\bfseries \hyperlink{clase6.10}{6~Modelo Biba (integridad)}}
  \item[]\hspace{1.2em}{\hyperlink{clase6.11}{6.1~Variante 1: Biba estricto}}
  \item[]\hspace{1.2em}{\hyperlink{clase6.12}{6.2~Variante 2: Biba Ring-Policy}}
  \item[]\hspace{1.2em}{\hyperlink{clase6.12}{6.3~Variante 3: Biba Low-Watermark}}
  \item[]{\bfseries \hyperlink{clase6.12}{7~Modelo Clark--Wilson (integridad comercial)}}
  \item[]\hspace{1.2em}{\hyperlink{clase6.13}{7.1~Reglas formales (clase práctica)}}
  \item[]{\bfseries \hyperlink{clase6.14}{8~Más allá de los modelos teóricos}}
  \item[]{\bfseries \hyperlink{clase6.15}{9~Mecanismos de control de acceso: MAC vs.\ DAC}}
  \item[]{\bfseries \hyperlink{clase6.16}{10~Control discrecional por extensión}}
  \item[]\hspace{1.2em}{\hyperlink{clase6.16}{10.1~Matriz de control de accesos}}
  \item[]\hspace{1.2em}{\hyperlink{clase6.16}{10.2~Listas de control de acceso (ACL)}}
  \item[]\hspace{1.2em}{\hyperlink{clase6.17}{10.3~Usuarios privilegiados, transferencias y revocaciones}}
  \item[]{\bfseries \hyperlink{clase6.18}{11~Control discrecional por comprensión}}
  \item[]\hspace{1.2em}{\hyperlink{clase6.18}{11.1~RBAC — control basado en roles}}
  \item[]\hspace{1.2em}{\hyperlink{clase6.18}{11.2~ABAC — control basado en atributos}}
  \item[]\hspace{1.2em}{\hyperlink{clase6.19}{11.3~RuBAC — control basado en reglas}}
  \item[]{\bfseries \hyperlink{clase6.19}{12~Resolución de múltiples reglas}}
  \item[]{\bfseries \hyperlink{clase6.20}{13~Reglas basadas en contexto}}
  \item[]\hspace{1.2em}{\hyperlink{clase6.20}{13.1~Zero-Trust y Open Policy Agent (OPA)}}
  \item[]{\bfseries \hyperlink{clase6.21}{14~Ejemplos reales de control de acceso}}
  \item[]\hspace{1.2em}{\hyperlink{clase6.21}{14.1~Unix / Linux file system}}
  \item[]\hspace{1.2em}{\hyperlink{clase6.22}{14.2~SELinux (Security-Enhanced Linux)}}
  \item[]\hspace{1.2em}{\hyperlink{clase6.22}{14.3~Windows file system}}
  \item[]\hspace{1.2em}{\hyperlink{clase6.22}{14.4~PostgreSQL: roles y policies}}
  \item[]\hspace{1.2em}{\hyperlink{clase6.23}{14.5~Squid Web Proxy}}
  \item[]\hspace{1.2em}{\hyperlink{clase6.23}{14.6~AWS (IAM)}}
  \item[]\hspace{1.2em}{\hyperlink{clase6.23}{14.7~API gateway}}
  \item[]{\bfseries \hyperlink{clase6.23}{15~OAuth 2.0 (Open Authorization)}}
  \item[]\hspace{1.2em}{\hyperlink{clase6.23}{15.1~El flujo de autorización}}
  \item[]\hspace{1.2em}{\hyperlink{clase6.24}{15.2~Refresh Token}}
  \item[]\hspace{1.2em}{\hyperlink{clase6.24}{15.3~Scope (alcance) — un mecanismo DAC}}
  \item[]{\bfseries \hyperlink{clase6.25}{~Lectura recomendada}}
\end{itemize}

\subsection*{\hyperref[src:clase7]{Clase 7 — Autenticación}}
\begin{itemize}[leftmargin=1.4em,itemsep=0pt,topsep=2pt]
  \item[]{\bfseries \hyperlink{clase7.3}{1~Introducción y contexto}}
  \item[]{\bfseries \hyperlink{clase7.3}{2~Conceptos fundamentales}}
  \item[]\hspace{1.2em}{\hyperlink{clase7.3}{2.1~Autenticación Autorización}}
  \item[]\hspace{1.2em}{\hyperlink{clase7.4}{2.2~El proceso de autenticación}}
  \item[]{\bfseries \hyperlink{clase7.4}{3~Sistema de autenticación}}
  \item[]\hspace{1.2em}{\hyperlink{clase7.5}{3.1~Ejemplo: autenticación con contraseñas}}
  \item[]{\bfseries \hyperlink{clase7.5}{4~Información de autenticación: los factores}}
  \item[]\hspace{1.2em}{\hyperlink{clase7.5}{4.1~Algo que conozco (conocimiento compartido)}}
  \item[]\hspace{1.2em}{\hyperlink{clase7.6}{4.2~Algo que tengo (posesión física)}}
  \item[]\hspace{1.2em}{\hyperlink{clase7.6}{4.3~Algo que soy (biométrico / inherente)}}
  \item[]\hspace{1.2em}{\hyperlink{clase7.6}{4.4~Contexto}}
  \item[]{\bfseries \hyperlink{clase7.6}{5~Claves (contraseñas)}}
  \item[]\hspace{1.2em}{\hyperlink{clase7.7}{5.1~Tipos de claves}}
  \item[]\hspace{1.2em}{\hyperlink{clase7.7}{5.2~Almacenamiento de claves}}
  \item[]{\bfseries \hyperlink{clase7.8}{6~Ataques a un sistema de autenticación}}
  \item[]\hspace{1.2em}{\hyperlink{clase7.8}{6.1~Adivinando claves: offline vs.\ online}}
  \item[]\hspace{1.2em}{\hyperlink{clase7.8}{6.2~Cómo se ``mejoran'' los ataques}}
  \item[]{\bfseries \hyperlink{clase7.9}{7~Prevenciones}}
  \item[]\hspace{1.2em}{\hyperlink{clase7.9}{7.1~Prevenciones generales}}
  \item[]\hspace{1.2em}{\hyperlink{clase7.9}{7.2~Subir la complejidad: fórmula de Anderson}}
  \item[]\hspace{1.2em}{\hyperlink{clase7.10}{7.3~Selección de claves — alternativas}}
  \item[]\hspace{1.2em}{\hyperlink{clase7.10}{7.4~Revisión proactiva de claves}}
  \item[]\hspace{1.2em}{\hyperlink{clase7.10}{7.5~Expiración de claves}}
  \item[]{\bfseries \hyperlink{clase7.10}{8~Almacenamiento seguro: Salting y PBKDF2}}
  \item[]\hspace{1.2em}{\hyperlink{clase7.10}{8.1~Salting}}
  \item[]\hspace{1.2em}{\hyperlink{clase7.11}{8.2~PBKDF2 (PKCS \#5 v2.0)}}
  \item[]{\bfseries \hyperlink{clase7.11}{9~Protocolos sin transmitir la clave}}
  \item[]\hspace{1.2em}{\hyperlink{clase7.12}{9.1~Challenge--Response (pregunta--respuesta)}}
  \item[]\hspace{1.2em}{\hyperlink{clase7.12}{9.2~EKE — Encrypted Key Exchange}}
  \item[]{\bfseries \hyperlink{clase7.12}{10~OTP — Claves de un solo uso}}
  \item[]\hspace{1.2em}{\hyperlink{clase7.13}{10.1~HOTP — RFC 4226 (basado en HMAC)}}
  \item[]\hspace{1.2em}{\hyperlink{clase7.13}{10.2~TOTP — RFC 6238 (basado en tiempo)}}
  \item[]{\bfseries \hyperlink{clase7.13}{11~Autenticación multifactor (MFA)}}
  \item[]{\bfseries \hyperlink{clase7.14}{12~Autenticación remota y federada (SSO)}}
  \item[]{\bfseries \hyperlink{clase7.14}{~Lectura recomendada}}
\end{itemize}

\subsection*{\hyperref[src:clase8]{Clase 8 — Principios de Diseño y Vulnerabilidades}}
\begin{itemize}[leftmargin=1.4em,itemsep=0pt,topsep=2pt]
  \item[]{\bfseries \hyperlink{clase8.3}{1~Introducción y contexto}}
  \item[]{\bfseries \hyperlink{clase8.3}{2~Principios de Diseño Seguro}}
  \item[]\hspace{1.2em}{\hyperlink{clase8.4}{2.1~El balance: tecnología, procesos y personas}}
  \item[]\hspace{1.2em}{\hyperlink{clase8.4}{2.2~Los 8 principios de Saltzer \& Schroeder}}
  \item[]{\bfseries \hyperlink{clase8.8}{3~Flujo de Información}}
  \item[]\hspace{1.2em}{\hyperlink{clase8.9}{3.1~Entropía H(x) e incertidumbre}}
  \item[]\hspace{1.2em}{\hyperlink{clase8.9}{3.2~Fórmulas de entropía}}
  \item[]\hspace{1.2em}{\hyperlink{clase8.10}{3.3~Ejemplo resuelto: los dos dados}}
  \item[]\hspace{1.2em}{\hyperlink{clase8.10}{3.4~El secreto compartido de Shamir, en términos de entropía}}
  \item[]\hspace{1.2em}{\hyperlink{clase8.11}{3.5~Flujo directo vs.\ indirecto}}
  \item[]{\bfseries \hyperlink{clase8.11}{4~Canales Ocultos y Aislamiento}}
  \item[]\hspace{1.2em}{\hyperlink{clase8.12}{4.1~Aislamiento: la solución ideal imposible}}
  \item[]\hspace{1.2em}{\hyperlink{clase8.12}{4.2~Máquinas virtuales y sandboxes}}
  \item[]{\bfseries \hyperlink{clase8.13}{5~Amenazas, Vulnerabilidades y Riesgo}}
  \item[]\hspace{1.2em}{\hyperlink{clase8.14}{5.1~Taxonomías de amenazas}}
  \item[]\hspace{1.2em}{\hyperlink{clase8.14}{5.2~Catálogo de vulnerabilidades}}
  \item[]\hspace{1.2em}{\hyperlink{clase8.16}{5.3~MITRE: CVE y CWE}}
  \item[]\hspace{1.2em}{\hyperlink{clase8.17}{5.4~Zero-day y el mercado de exploits}}
  \item[]\hspace{1.2em}{\hyperlink{clase8.17}{5.5~OWASP Top 10}}
  \item[]{\bfseries \hyperlink{clase8.18}{6~Modelado de Amenazas (Threat Modeling) y STRIDE}}
  \item[]\hspace{1.2em}{\hyperlink{clase8.19}{6.1~Modelo STRIDE}}
  \item[]\hspace{1.2em}{\hyperlink{clase8.19}{6.2~Data Flow Model (DFD)}}
  \item[]{\bfseries \hyperlink{clase8.20}{~Lectura recomendada}}
\end{itemize}

\subsection*{\hyperref[src:clase9]{Clase 9 — Flujo de Información y Malware}}
\begin{itemize}[leftmargin=1.4em,itemsep=0pt,topsep=2pt]
  \item[]{\bfseries \hyperlink{clase9.3}{1~Introducción y contexto}}
  \item[]{\bfseries \hyperlink{clase9.3}{2~El problema: control de acceso vs.\ flujo de información}}
  \item[]{\bfseries \hyperlink{clase9.4}{3~El problema del confinamiento}}
  \item[]\hspace{1.2em}{\hyperlink{clase9.4}{3.1~Modelo ideal: aislación total}}
  \item[]{\bfseries \hyperlink{clase9.5}{4~Canales encubiertos (Covert Channels)}}
  \item[]\hspace{1.2em}{\hyperlink{clase9.5}{4.1~Canales de temporización (Timing Channels)}}
  \item[]\hspace{1.2em}{\hyperlink{clase9.6}{4.2~Canales de almacenamiento (Storage Channels)}}
  \item[]\hspace{1.2em}{\hyperlink{clase9.6}{4.3~Ejemplos de canales encubiertos}}
  \item[]{\bfseries \hyperlink{clase9.6}{5~Ataques de canal lateral (Side-Channel Attacks)}}
  \item[]\hspace{1.2em}{\hyperlink{clase9.7}{5.1~Ejemplos}}
  \item[]{\bfseries \hyperlink{clase9.7}{6~Control de flujo}}
  \item[]\hspace{1.2em}{\hyperlink{clase9.8}{6.1~Control de flujo en tiempo de compilación}}
  \item[]\hspace{1.2em}{\hyperlink{clase9.10}{6.2~Control de flujo en tiempo de ejecución: aislamiento}}
  \item[]{\bfseries \hyperlink{clase9.11}{7~Malware}}
  \item[]\hspace{1.2em}{\hyperlink{clase9.12}{7.1~Virus / Worm}}
  \item[]\hspace{1.2em}{\hyperlink{clase9.12}{7.2~Troyano}}
  \item[]\hspace{1.2em}{\hyperlink{clase9.12}{7.3~Rootkit}}
  \item[]\hspace{1.2em}{\hyperlink{clase9.13}{7.4~Spyware}}
  \item[]\hspace{1.2em}{\hyperlink{clase9.13}{7.5~Ransomware}}
  \item[]\hspace{1.2em}{\hyperlink{clase9.13}{7.6~Botnet}}
  \item[]\hspace{1.2em}{\hyperlink{clase9.13}{7.7~Ciclo de ataque del malware}}
  \item[]{\bfseries \hyperlink{clase9.15}{8~Práctica: CWE Top 25 — Errores de software más peligrosos}}
  \item[]\hspace{1.2em}{\hyperlink{clase9.15}{8.1~Metodología del ranking}}
  \item[]\hspace{1.2em}{\hyperlink{clase9.15}{8.2~Evolución de la clasificación}}
  \item[]\hspace{1.2em}{\hyperlink{clase9.15}{8.3~El CWE Top 25 — 2025, agrupado por categoría}}
  \item[]\hspace{1.2em}{\hyperlink{clase9.18}{8.4~Debilidades recurrentes (entran y salen del ranking)}}
  \item[]\hspace{1.2em}{\hyperlink{clase9.18}{8.5~Ejemplo: impacto y fase de un buffer overflow}}
  \item[]{\bfseries \hyperlink{clase9.19}{~Lectura recomendada}}
\end{itemize}

\subsection*{\hyperref[src:clase10]{Clase 10 — Seguridad en la Empresa}}
\begin{itemize}[leftmargin=1.4em,itemsep=0pt,topsep=2pt]
  \item[]{\bfseries \hyperlink{clase10.2}{1~Seguridad en la Empresa}}
  \item[]\hspace{1.2em}{\hyperlink{clase10.2}{1.1~Gobernanza (Governance)}}
  \item[]\hspace{1.2em}{\hyperlink{clase10.2}{1.2~Gestión de riesgos (Risk Management)}}
  \item[]\hspace{1.2em}{\hyperlink{clase10.4}{1.3~Políticas de seguridad}}
  \item[]\hspace{1.2em}{\hyperlink{clase10.4}{1.4~Estrategias: Defense-in-depth vs.\ Zero Trust}}
  \item[]\hspace{1.2em}{\hyperlink{clase10.8}{1.5~Privileged Access Management (PAM)}}
  \item[]\hspace{1.2em}{\hyperlink{clase10.8}{1.6~Prevención de vulnerabilidades: DevSecOps}}
  \item[]\hspace{1.2em}{\hyperlink{clase10.10}{1.7~Protección de la nube}}
  \item[]\hspace{1.2em}{\hyperlink{clase10.10}{1.8~Evaluaciones de seguridad}}
\end{itemize}

\subsection*{\hyperref[src:clase10b]{Clase 10b — Pentesting}}
\begin{itemize}[leftmargin=1.4em,itemsep=0pt,topsep=2pt]
  \item[]{\bfseries \hyperlink{clase10b.2}{1~Pentesting}}
  \item[]\hspace{1.2em}{\hyperlink{clase10b.2}{1.1~Formas de verificar un sistema: verificación formal vs.\ pentesting}}
  \item[]\hspace{1.2em}{\hyperlink{clase10b.3}{1.2~Objetivos y marco ético (ethical hacking)}}
  \item[]\hspace{1.2em}{\hyperlink{clase10b.4}{1.3~Metodología informal (preparación)}}
  \item[]\hspace{1.2em}{\hyperlink{clase10b.4}{1.4~Metodología de Hipótesis de Falla (LOS PASOS)}}
  \item[]\hspace{1.2em}{\hyperlink{clase10b.8}{1.5~Herramientas y técnicas mencionadas}}
  \item[]\hspace{1.2em}{\hyperlink{clase10b.8}{1.6~Ejemplos}}
  \item[]\hspace{1.2em}{\hyperlink{clase10b.10}{1.7~Limitaciones: el pentesting NO garantiza la seguridad}}
  \item[]\hspace{1.2em}{\hyperlink{clase10b.10}{1.8~Lectura recomendada}}
\end{itemize}

\subsection*{\hyperref[src:clase11]{Clase 11 — Protección de Datos}}
\begin{itemize}[leftmargin=1.4em,itemsep=0pt,topsep=2pt]
  \item[]{\bfseries \hyperlink{clase11.2}{1~Protección de Datos}}
  \item[]\hspace{1.2em}{\hyperlink{clase11.2}{1.1~El problema y el ciclo de protección}}
  \item[]\hspace{1.2em}{\hyperlink{clase11.3}{1.2~Marco legal y regulaciones}}
  \item[]\hspace{1.2em}{\hyperlink{clase11.5}{1.3~Estados del dato y encripción}}
  \item[]\hspace{1.2em}{\hyperlink{clase11.6}{1.4~Cumplimiento, detección y respuesta}}
  \item[]\hspace{1.2em}{\hyperlink{clase11.7}{1.5~Encripción Homomórfica}}
\end{itemize}
